% ============================================================
%  Глава 10 — Двухфазный сценарий TURN → DRIVE
% ============================================================
\chapter{Двухфазный сценарий: TURN, затем DRIVE}
\label{ch:turn_phase}

\section{Вопрос}

В прошлой главе сценарий ехал прямо «вперёд от текущей ориентации».
А если пользователь \emph{выбрал} цель не прямо перед роботом, а
под углом $\varphi$ (например, $+45^\circ$ влево)? Что делать?

\begin{ideabox}
Можно было бы дать MPC сразу полную задачу: «доехать до точки $(D\cos\varphi,
D\sin\varphi)$, развернувшись по ходу». Но это плохая идея для линеаризованной
вокруг \emph{прямолинейного} движения модели: $\varphi = 45^\circ$ ---
далеко не «малое отклонение», линейная модель там врёт.

\textbf{Решение --- двухфазный сценарий:}
\begin{enumerate}
  \item \textbf{TURN.} Робот стоит на месте, поворачивается к курсу
        $\varphi$. После завершения --- курс точно $\varphi$, позиция
        не изменилась.
  \item \textbf{DRIVE.} Робот едет $D$ метров, удерживая курс $\varphi$.
        Это уже та же задача из главы~\ref{ch:scenario}, только
        $\theta_\text{ref} = \varphi$ вместо $0$.
\end{enumerate}
Между фазами --- мгновенный переход (без паузы), но \emph{с разными
ограничениями} на $u_\omega$ и \emph{другим} reference.
\end{ideabox}

\section{Фаза TURN}

\subsection{Reference}

В фазе TURN мы хотим, чтобы робот стоял ($s = 0$, $v = 0$) и развернулся
к курсу $\varphi$:

\begin{formulabox}[title={Reference фазы TURN}]
\begin{equation}
  \vect{x}^\text{ref}_\text{turn} = \begin{pmatrix} 0\\ 0\\ \varphi\\ 0\\ 0 \end{pmatrix}.
  \label{eq:turn_ref}
\end{equation}
\end{formulabox}

Подадим этот reference в \texttt{mpc.step(x, x\_ref)}. Полученное
$\vect{u}_\text{mpc}$ будет «оптимальным» --- но MPC не знает, что мы
хотим \emph{чистый разворот}, поэтому без дополнительных охранников
$u_v$ окажется не нулевым (особенно если $\varphi$ большой).

\subsection{Принудительные ограничения фазы}

В коде это решают \textbf{пост-обработкой}:

\begin{codebox}[title={\filepath{pi\_nodes/nodes/mps\_node.py:384}, фаза TURN}]
\begin{lstlisting}[language=Python]
def _tick_turn(self, run, x):
    phi = run.target_heading

    # Курс совпал с целью → переход в DRIVE.
    if abs(_normalize_angle(x[_THETA] - phi)) < self._turn_tol:
        run.phase = 'drive'
        run.drive_t = 0.0
        return True

    x_ref = np.array([0.0, 0.0, phi, 0.0, 0.0])
    u = self._mpc.step(x, x_ref=x_ref)

    # Чистое вращение: ход — в ноль; ω — свой (более высокий) кап.
    u[0] = 0.0
    u[1] = max(-self._omega_max_turn,
               min(self._omega_max_turn, u[1]))

    self._publish_cmd_and_telemetry(run, x, u)
    # ... watchdog, timeout
\end{lstlisting}
\end{codebox}

\textbf{Что происходит:}
\begin{enumerate}
  \item \texttt{u[0] = 0.0} --- принудительный ход в ноль. MPC мог
        захотеть двигаться, но мы не разрешаем.
  \item \texttt{u[1]} --- угловая скорость зажимается в
        $[-\omega_\text{max,turn}, +\omega_\text{max,turn}]$, где
        $\omega_\text{max,turn} = 1.0$~рад/с --- \emph{больше}, чем в
        DRIVE-фазе ($0.5$). На месте можно крутиться быстрее.
\end{enumerate}

\begin{notebox}
Почему \emph{разный} cap на $\omega$ в TURN и DRIVE? Потому что:
\begin{itemize}
  \item В TURN мы хотим быстро развернуться --- иначе зачем вообще
        фаза.
  \item В DRIVE \emph{большой} $\omega$ при движении --- это «винт»,
        робот съезжает с курса. Лучше держать узкий cap, чтобы
        регулятор был «нежным» в коррекциях.
\end{itemize}
Два параметра: \texttt{mps.scenario.omega\_max\_in\_turn = 1.0} и
\texttt{omega\_max\_in\_forward = 0.5}.
\end{notebox}

\subsection{Условие завершения TURN}

Переходим в DRIVE, когда курс совпал с целью с заданным допуском:

\begin{equation}
  \big|\,\text{wrap}_{[-\pi, \pi]}(\theta - \varphi)\,\big| < \varepsilon_\text{turn}.
\end{equation}

В коде $\varepsilon_\text{turn} = 0.05$~рад $\approx 2.9^\circ$ ---
поправка \texttt{mps.scenario.turn\_tolerance\_rad}.

\begin{warnbox}
\textbf{Wrap-around в $[-\pi, \pi]$ --- важно!} Если $\theta = +179^\circ$
и $\varphi = -179^\circ$, разность «по-наивному» $= 358^\circ$ --- огромная,
хотя физически они почти совпадают.

Поэтому используем функцию \texttt{\_normalize\_angle(a)}:
\begin{equation}
  \text{wrap}_{[-\pi, \pi]}(a) = ((a + \pi) \bmod 2\pi) - \pi.
\end{equation}
В коде:
\begin{lstlisting}[language=Python]
def _normalize_angle(a: float) -> float:
    return (a + math.pi) % (2 * math.pi) - math.pi
\end{lstlisting}

Без этого, в редких случаях, MPC получал бы reference-ошибку курса
$\Delta\theta \sim 2\pi$ --- умножал бы её на $\mat{K}$ --- получал бы
безумно большое $u_\omega$ --- упирался в насыщение и крутился в
обратную сторону.
\end{warnbox}

\subsection{Тайм-аут фазы TURN}

Чтобы при поломанной IMU TURN не зациклился навсегда, есть собственный
тайм-аут: $t_\text{turn} \le 10$~с (параметр
\texttt{mps.scenario.turn\_timeout\_s}). Не успели --- завершаем
прогон с \texttt{status='timeout'}.

\section{Переход TURN \texorpdfstring{$\to$}{->} DRIVE}

Когда условие завершения TURN сработало, \texttt{\_tick\_turn}
возвращает \texttt{True}, и в \texttt{\_tick} тем же тиком вызывается
\texttt{\_tick\_drive(run, x)}. Между ними:

\begin{itemize}
  \item не теряем тик (не пропускаем 20~мс на лоське);
  \item \texttt{run.phase = 'drive'} --- фаза переключилась;
  \item \texttt{run.drive\_t = 0.0} --- часы DRIVE-фазы стартуют с нуля,
        отдельно от общего $t$ (для ramp $s_\text{ref}$ и DRIVE-тайм-аута).
\end{itemize}

\begin{codebox}[title={\filepath{mps\_node.py:366}, общий tick}]
\begin{lstlisting}[language=Python]
def _tick(self):
    with self._lock:
        run = self._run
        if run is None or self._fsm_state != 'DRIVE_FORWARD_MPS':
            return
        x = self._x_meas.copy()
    # Относительные координаты.
    x[_S] = x[_S] - run.s_start
    x[_THETA] = _normalize_angle(x[_THETA] - run.theta_start)

    # Двухфазный сценарий.
    if run.phase == 'turn':
        if not self._tick_turn(run, x):
            return  # ещё крутимся / прогон завершён
        # TURN завершился ЭТИМ тиком → продолжаем в DRIVE
    self._tick_drive(run, x)
\end{lstlisting}
\end{codebox}

\section{Фаза DRIVE}

После переключения в DRIVE логика та же, что в главе~\ref{ch:scenario},
\emph{с одной поправкой}: курс reference --- $\varphi$, а не $0$.

\[
  \vect{x}^\text{ref}_\text{drive}(k) =
  \begin{pmatrix} s_\text{ref}(k)\\ v_\text{target}\\ \varphi\\ 0\\ 0 \end{pmatrix},
  \qquad
  s_\text{ref}(k) = \min(D, k\,\Delta t \cdot v_\text{target}).
\]

\begin{ideabox}
\textbf{Почему именно так, а не «следовать дуге»?} Робот в DRIVE
\emph{уже} ориентирован к $\varphi$ (фаза TURN об этом позаботилась).
Он едет вперёд в своей системе координат, а курс держит постоянным.
Это самая простая постановка, и она работает.

Если бы мы хотели ехать по \emph{дуге} (например, для плавного объезда
препятствия), то и reference, и линеаризация были бы сложнее: вместо
постоянной опорной $(v_0, 0)\T$ надо было бы линеаризовать вдоль дуги
с $\bar\omega \ne 0$. Это уже не учебный сценарий, а полноценная
trajectory tracking.
\end{ideabox}

\section{Timeline прогона}

\begin{center}
\begin{tikzpicture}[
  scale=1.0, font=\small,
  bar/.style={very thick, line cap=round},
]
  % Ось времени
  \draw[-{Stealth}] (-0.2, 0) -- (10.5, 0) node[right]{$t$, с};
  % Метки
  \foreach \x in {0,2,4,6,8,10}
    \draw (\x, -0.05) -- (\x, 0.05) node[above, font=\tiny]{\x};
  % TURN полоска
  \draw[bar, red!70!black]  (0, -0.6) -- (2.5, -0.6);
  \node[red!70!black] at (1.25, -0.95) {\textbf{TURN}};
  \node[font=\tiny, red!70!black] at (1.25, -1.25) {$\theta \to \varphi$};
  % DRIVE полоска
  \draw[bar, blue!70!black] (2.5, -0.6) -- (9.5, -0.6);
  \node[blue!70!black] at (6, -0.95) {\textbf{DRIVE}};
  \node[font=\tiny, blue!70!black] at (6, -1.25) {$s \to D$, $\theta = \varphi$};
  % Переход
  \draw[dashed, gray] (2.5, 0.3) -- (2.5, -1.5);
  \node[font=\scriptsize, gray] at (2.5, 0.55) {TURN$\to$DRIVE};
  % FINISHED
  \draw[bar, green!50!black] (9.5, -0.6) -- (10.0, -0.6);
  \node[green!50!black, font=\tiny] at (10.4, -0.95) {STOP};
\end{tikzpicture}
\end{center}

\textbf{Поля \texttt{\_RunState}, обслуживающие двухфазность:}

\begin{itemize}
  \item \texttt{run.t} --- монотонное \emph{суммарное} время от старта
        прогона. Растёт всегда. Используется для тайм-аута фазы TURN
        и как абсцисса телеметрии.
  \item \texttt{run.drive\_t} --- часы только DRIVE-фазы, стартуют с
        $0$ в момент переключения. Используются для ramp $s_\text{ref}$
        и DRIVE-тайм-аута.
  \item \texttt{run.phase} --- \texttt{'turn'} либо \texttt{'drive'}.
\end{itemize}

\section{Граничные случаи}

\begin{whatifbox}
\textbf{$\varphi = 0$.} Условие завершения TURN срабатывает на первом
же тике (если курс точно $0$): $\theta_\text{start} = 0 \Rightarrow x[\theta]
= 0 \Rightarrow |0 - 0| = 0 < \varepsilon_\text{turn}$. Сценарий
\emph{скипает} фазу TURN и идёт сразу в DRIVE. Это и есть «прямо вперёд»
из предыдущей главы.

\textbf{$\varphi = \pi$ (поворот назад).} Худший случай: робот должен
развернуться на $180^\circ$. Через wrap-around $\theta - \varphi$ всё
ещё корректно (нормализуется в $[-\pi, \pi]$), но при $\theta = 0$,
$\varphi = \pi$ ошибка ровно $\pm\pi$ --- знак $u_\omega$
неопределён. На практике первый тик MPC выберет какое-то направление,
робот начнёт крутиться, и дальше всё разрешится. Скорее всего проект
зафиксирует $\varphi \in (-\pi, \pi]$, не \emph{строгий} $\pi$, но валидация
в \texttt{\_on\_scenario\_run} допускает.

\textbf{Watchdog в TURN.} Если на каком-то тике нет одометрии $> 3$
тиков подряд, мы не можем достоверно судить, повернулся ли робот.
TURN завершается \texttt{error}, FSM возвращается в IDLE.
\end{whatifbox}

\section{Зачем это всё было нужно}

\begin{ideabox}
Думаем учебно: какие были \emph{альтернативы}?

\textbf{(А) Не разделять фазы, дать MPC сразу полную задачу с целевой
точкой $(D\cos\varphi, D\sin\varphi)$.} Тогда модель должна была бы
быть \emph{инвариантной} к ориентации --- то есть линеаризация вокруг
прямолинейного движения уже не работает; нужна нелинейная модель или
linearization gain scheduling.

\textbf{(Б) Сначала повернуться, потом ехать --- но через два разных
регулятора.} TURN --- простой П-регулятор на угле, DRIVE --- MPC.
Так у нас и было раньше (старый PID). Но это значит, что в коде
\emph{два} регулятора, и переключение между ними может давать всплеск
по \texttt{cmd\_vel}. Также --- два разных набора тюнинга, два моделирования.

\textbf{(В, выбранный) Один и тот же MPC, разные reference и капы
управления в каждой фазе.} Минусы: дополнительная сложность
переключения. Плюсы: один контроллер --- одна модель --- единое поведение.
Прозрачно для учебного эксперимента: те же $\mat{Q}, \mat{R}$ работают
в обеих фазах.
\end{ideabox}

\section{Что мы получили}

\begin{itemize}
  \item Сценарий стал \textbf{двухфазным}: TURN $\to$ DRIVE.
  \item Используется тот же MPC, но с разными reference и капами:
        TURN режим --- $u_v = 0$, $|u_\omega| \le 1.0$;
        DRIVE режим --- $u_v$ свободно, $|u_\omega| \le 0.5$.
  \item Переход --- мгновенный, тем же тиком, без пропуска
        управления.
  \item Угол $\varphi$ работает корректно через wrap-around в $[-\pi, \pi]$.
  \item Граничные случаи ($\varphi = 0$, watchdog) обрабатываются
        самой логикой.
\end{itemize}

Остался один вопрос: \textbf{откуда берётся $\varphi$}? Это
3D-пикер цели --- последняя интересная математика в этом учебнике.
